\section{Introduction}\label{sec:introduction}

Let \(A_n\) be a symmetric \(n\times n\) random matrix whose entries
on and above the diagonal are independent and uniform on \(\{-1,1\}\).
We study the probability that \(A_n\) is singular. A specified pair of
rows is equal with probability \(2^{-n}\): the \(n-2\) comparisons away
from the two diagonal positions and the two remaining diagonal
conditions are independent. Thus singularity already occurs at this
exponential scale.

\begin{theorem}\label{thm:main}
For the symmetric Bernoulli matrix \(A_n\),
\[
 \Pr(\det A_n=0)=\left(\frac12+o(1)\right)^n
                =2^{-n+o(n)}
 \qquad\text{as }n\to\infty.
\]
\end{theorem}

The corresponding problem for a matrix with all \(n^2\) entries
independent has a long history. Koml\'os proved that a matrix with
independent uniform \(0/1\) entries is nonsingular with probability
tending to one \cite{Komlos}. Tikhomirov proved the sharp rate
\((1/2+o(1))^n\) for independent uniform signs
\cite{Tikhomirov}. Symmetry changes the exposure of the matrix:
revealing a row also reveals entries in all later rows. Arguments
based on independent rows therefore need a different source of
anti-concentration.

Costello, Tao, and Vu proved that a random symmetric Bernoulli matrix
is nonsingular with probability tending to one, obtaining the bound
\(O(n^{-1/8+\delta})\) for every fixed \(\delta>0\)
\cite{CTV}. Their proof developed quadratic Littlewood--Offord
estimates to handle the dependence between a row and its matching
column. Nguyen improved the bound to \(O_C(n^{-C})\) for every fixed
\(C>0\) \cite{Nguyen}, and Vershynin obtained a bound of the form
\(\exp(-n^c)\) for the sign model, as a special case of his theorem
for a class of subgaussian ensembles \cite{Vershynin}. Further
quantitative advances for this model gave bounds
of the forms
\(\exp(-c n^{1/4}\sqrt{\log n})\) \cite{FerberJain},
\(\exp(-c\sqrt n)\) \cite{CamposMattosMorrisMorrison},
\(\exp(-c\sqrt n(\log n)^{1/4})\) \cite{JainSahSawhney}, and
\(\exp(-c\sqrt{n\log n})\) \cite{CJMSRevisited}.
Campos, Jenssen, Michelen, and Sahasrabudhe proved an exponential
bound \(\exp(-cn)\), using inverse Littlewood--Offord estimates for
conditioned random walks and the method of inversion of randomness
\cite{SymmetricExponential}. Related work gives
quantitative control of the least singular value
\cite{SymmetricLeastSingular} and of larger rank deficiencies
\cite{HanRank}. Theorem~\ref{thm:main} identifies the leading
exponential rate for the specified sign model.

Our proof begins with an exact reduction to symmetric matrices of
bits. A Schur complement converts the sign matrix in dimension \(n\)
to a binary matrix in dimension \(N=n-1\). A reduction to corank one
then turns singularity into a question about appending a new column to
a nonsingular matrix \(B\). The relevant bit columns \(z\) satisfy
both \(z^{\mathsf T}B^{-1}z\in\bits\) and
\(\|B^{-1}z\|_\infty\le1\). It suffices to prove that their expected
number is \(2^{o(N)}\).

For a typical matrix \(B\), the column count is controlled by an
arithmetic potential \(\Psi(B)\in[0,1]\), defined in
Section~\ref{sec:framework}. The discriminant group
\(G_B=\Z^m/B\Z^m\) carries the symmetric pairing
\[
 ([x],[y])\longmapsto x^{\mathsf T}B^{-1}y\pmod{\Z}.
\]
Admissible classes have zero self-pairing. We partition them so that
pairings within each part have controlled denominators, with the cost
expressed in terms of the invariant factors. A robust subset of a large
part yields a lattice of controlled covolume on the graph of \(B^{-1}\).
A cover for rational subspaces of controlled height makes rank estimates
simultaneous for the span selected from these columns. Comparing a lower
determinant bound from the pairing with an upper bound from typical rank
and spectral properties proves the column estimate. The reusable lattice
arguments control covolumes and count short vectors after removing a
determinant-minimizing sublattice; the Reverse Minkowski Theorem of Regev
and Stephens-Davidowitz~\cite{ReverseMinkowski} is applied to its
orthogonal quotient.

For a nonsingular extension of width one, the old and new discriminant
groups have a common quotient. A kernel count modulo prime powers uses
symmetry to restrict the possible layer changes, while retaining
prescribed layers forces the new column into an integer lattice. A
second application of the lattice tools bounds cube intersections
simultaneously for all such lattices; it compares the cost of specifying
bases of cube differences with the number of matrix entries they
constrain. Counting partial certificates of retained layers and
optimizing their size then pays for the increase of \(\Psi\) in a
conditional exponential moment. For extensions of width two, a gate
requires both new columns to be admissible, and the static column
estimate directly pays for the probability of this gate.

Finally, every nonsingular matrix admits a path of nonsingular principal
minors with steps of width one or two, with the width-two steps satisfying
the gate. A separate elementary estimate for retaining high layers,
together with a bound on coranks modulo all primes, uses a long initial
part of the path to make a large high-layer tail unlikely anywhere in
a short terminal window. Outside this exceptional event, if the potential
is small at a node in that
window, the local conditional estimate controls the terminal moment.
Otherwise a bounded auxiliary statistic decreases at every step except
those with substantial retention or an expensive width-two gate. Many
such steps are necessary, and their conditional probability bounds make
this alternative negligible. The long initial part controls the
high-layer tail, while the short window limits the accumulation of local
errors.
